\section{Related Work}\label{sec:relatedwork}

\subsection{Characterizing bugs and patches in APR
}
Past work has characterized bugs in popular datasets used in automated program repair and more broadly in open source repositories. Sobreira et al~\cite{sobreira2018dissection} characterized the patches for the popular Java-based program repair benchmark. Zhong et al~\cite{zhong2015empirical} explored bugs and their fixes in five large open source projects.  Similar to our GITS/SWE-Bench comparison, these works measured patch size and spread. In contrast to their work, we draw our bugs from GITS, reflecting real tasks for Google developers, and these are multilingual (covering Python, C++, Dart, Go, and Javascript/TypeScript).

\subsection{Agent-based APR}

SWE-Agent~\cite{sweagent}, which directly inspired the implementation of Passerine, implements a ReAct-style loop with access to a set of bash-style APIs as tools. In contrast to the fully dynamic approach present in SWE-Agent, RepairAgent~\cite{repairagent} presents an agent whose actions are constrained by a state machine, disallowing certain tool use sequences. AutoCodeRover~\cite{autocoderover} is an agent-based repair system with tools that exploit explicit program information, such as class/method-based definition search and test-based information for localization. SpecRover~\cite{specrover} augments AutoCodeRover by defining a natural-language specification for the expected behavior at each candidate repair location and introduces a patch review agent. CodeR~\cite{codeR} decomposes APR into multiple subagent tasks that are coordinated with a task graph created and reviewed by a “manager” agent. Similarly, AutoDev~\cite{autodev} employs a multi-agent approach to tackle software engineering tasks beyond program repair, such as test generation and general code completion. MarsCode Agent~\cite{marscode} blends a dynamic, iterative approach to program repair with a traditional generate-and-validate pipeline in a multi-agent repair framework. OpenDevin~\cite{opendevin} — an open-source agent platform inspired by the commercial Devin software engineering agent~\cite{devin} — can be used to build agent-based solutions for tasks in software engineering and other domains.

Similarly to this body of work, Passerine is an agent-based framework. However, our contribution is focused exclusively on program repair. Like many of these approaches, Passerine employs tools in a ReAct style and evaluates its performance using a test-suite-based metric. Our current implementation is aligned with the vision presented in SWE-Agent and exposes a subset of the commands available to that agent. However, in contrast to these systems, Passerine has been developed as a proof-of-concept to tackle bugs in Google’s development environment and as a result has been implemented to use custom tools typically used by Google developers. As a result, our contribution constitutes the first systematic study of  an agent-based APR system in a large industrial codebase.


\subsection{LLM-based APR
}

Prior work has also explored the use of LLMs for fixing and identifying bugs without the need for agents. AlphaRepair~\cite{alpharepair} showed that LLMs could be used to perform zero-shot program repair by framing fixing code as an in-filling task. Agentless~\cite{agentless} obtained results competitive with agent-based approaches on SWE-Bench by employing a fixed repair pipeline, consisting of a hierarchical fault localization approach and multi-patch generation paired with filtering and ranking of patches. ChatRepair~\cite{chatrepair} made use of a conversational-style repair, where the LLM incrementally gains information about patch candidates that fail tests and plausible patch candidates to generate more diverse and accurate patches.
Zhang et al~\cite{yang2024large} also showed that LLMs can be effectively fine-tuned and used as fault localizers for real bugs. 

Like these works, Passerine is also focused on APR, but in contrast it is an agent-based system. Futhermore, Passerine has been designed to tackle bugs within Google’s development environment. Our current evaluation does not use a model fine-tuned for repair nor do we provide meaningful few-shot examples, beyond fixed examples in command APIs to explain how a command is used.
